% === VI. CONCLUSION ===
\section{Conclusion}
\label{sec:conclusion}

We presented \system{}, a multimodal AI rehabilitation system for elderly users. The system integrates real-time whole-body 3D pose estimation via RTMW3D-L (133 keypoints covering the body, hands, and face), multi-indicator facial state detection with OpenFace~3.0, voice-interactive LLM coaching, and a clinically inspired six-dimensional scoring framework.

The key contributions are: (1)~to our knowledge, the first integrated rehabilitation system combining whole-body 3D pose estimation, multi-indicator facial state detection, voice-interactive LLM coaching, and clinically-inspired kinematic scoring in a single real-time framework; (2)~a whole-body 3D pose estimation pipeline using RTMW3D-L from MMPose with full 133-keypoint coverage running in real time on a single consumer GPU; (3)~a multi-indicator facial state detection pipeline integrating PSPI-inspired pain estimation, PERCLOS-based fatigue monitoring, blink/yawn detection, and engagement analysis; and (4)~a DTW-based scoring framework validated on two public rehabilitation benchmarks, achieving significant clinical correlation on KIMORE (Spearman $\rho = 0.347$, $p < 10^{-11}$) and near-perfect exercise discriminability on UCO (AUC = 0.974, 48-point separation).

This paper is a feasibility study: the pipeline is computationally practical, and the scoring framework shows significant clinical signal, but clinical validation with elderly participants, ground-truth pose-accuracy benchmarking, and per-dimension weight calibration against clinician ratings remain future work.

Future work includes: (1)~per-dimension scoring weight calibration against physiotherapist assessments; (2)~clinical user studies with elderly participants to assess usability and therapeutic effectiveness; (3)~metric-scale calibration to enable direct comparison with clinical goniometer measurements; (4)~edge-based LLM deployment to eliminate internet connectivity requirements; and (5)~multi-camera fusion to address self-occlusion limitations.
